% kubernetes-vulns.tex
% Thesis section: Extending MAD Goat to the Kubernetes Layer
% Compiled from thesis/kubernetes-vulns/ with:
%   pdflatex -shell-escape test-main.tex

\section{Extending the MAD Goat Benchmark to the Kubernetes Layer}
\label{sec:k8s-vulns}

\subsection{Motivation and Context}
\label{sec:k8s-motivation}

The MAD~Goat benchmark~\cite{mad-goat-repo} was originally designed to surface web-application and microservice-level vulnerabilities in containerised workloads.
As Kubernetes adoption has grown, the attack surface has shifted from individual containers to the orchestration layer itself: cluster RBAC, admission controllers, network policies, kubelet configuration, and audit infrastructure all introduce exploitable misconfigurations that are invisible to application-layer scanners~\cite{shamim-2020,bose-2023}.
The OWASP Kubernetes Top~10 2025 edition~\cite{owasp-k8s-2025} codifies ten categories of such weaknesses, updating and renumbering the 2022 list~\cite{owasp-k8s-2022} (Table~\ref{tab:owasp-crosswalk}).
This section extends MAD~Goat to the Kubernetes layer by injecting all ten categories into the benchmark's Helm charts and cluster bootstrap configuration, providing a reproducible, toggle-controlled environment for evaluating both active exploits and static analysis tools.

% Table 1: OWASP Kubernetes Top 10 crosswalk, 2022 vs. 2025
\begin{table}[ht]
\centering
\caption{OWASP Kubernetes Top~10 Crosswalk: 2022 vs.\ 2025~\cite{owasp-k8s-2025,owasp-k8s-2022}}
\label{tab:owasp-crosswalk}
\begin{tabular}{@{}clcl@{}}
\toprule
\textbf{2025} & \textbf{2025 Title} & \textbf{2022} & \textbf{2022 Title} \\
\midrule
K01 & Insecure Workload Configurations      & K01  & Insecure Workload Configurations \\
K02 & Overly Permissive RBAC                & K03  & Overly Permissive RBAC \\
K03 & Secrets Management Failures           & K08  & Secrets Management Failures \\
K04 & Lack of Centralised Policy Enforcement & K04  & Policy Enforcement \\
K05 & Inadequate Network Segmentation       & K07  & Network Segmentation Failures \\
K06 & Broken Authentication Mechanisms      & ---  & \textit{(new in 2025)} \\
K07 & Misconfigured Cluster Components      & K09  & Misconfigured Cluster Components \\
K08 & Supply Chain Vulnerabilities          & K02  & Supply Chain Vulnerabilities \\
K09 & Misconfigured Authentication Defaults & K06  & Broken Authentication \\
K10 & Inadequate Audit Logging              & K05  & Inadequate Logging \\
\bottomrule
\end{tabular}
\end{table}

\subsection{Methodology}
\label{sec:k8s-methodology}

Each vulnerability entry is controlled by a boolean flag in the chart's \texttt{values.yaml} (default: \texttt{true} = vulnerable).
Chart-level entries (\texttt{scope: chart}) live in \texttt{helm/madgoat} or \texttt{helm/madgoat-infra} and can be toggled with \texttt{helm upgrade}.
Cluster-level entries (\texttt{scope: cluster}) are applied via \texttt{kubectl apply -k k8s/insecure-defaults/} and require a full Minikube restart to remediate; the \texttt{make reset-safe} target automates this.
Kubelet flags (K07) are passed as \texttt{minikube start -{}-extra-config} arguments and are therefore not Helm-controlled.
Two entries (K08) are conceptual only: the IMDS credential-theft attack surface does not exist in a local Minikube cluster and is documented with references to cloud-vendor behaviour~\cite{aws-imds,azure-imds}.

For each \emph{full}-class entry the benchmark provides:
\begin{itemize}
  \item A \texttt{values.yaml} boolean flag that injects or suppresses the misconfiguration.
  \item A verification script in \texttt{scripts/exploits/} that exits~0 (PASS) when the vulnerable condition is confirmed and exits~1 (FAIL/safe) otherwise.
  \item Code snippets (Listings below) illustrating the vulnerable and remediated manifests.
\end{itemize}

The ground-truth inventory is maintained in \texttt{thesis/kubernetes-vulns/data/vulnerabilities.yaml}; Tables~\ref{tab:vuln-inventory} and~\ref{tab:scanner-coverage} are generated from it automatically by \texttt{scripts/gen-tables.py}~\cite{trivy-k8s,kics-tool,kubelinter,kubescape}.

% Generated by scripts/gen-tables.py — do not edit manually
\begin{longtable}{@{}lp{5cm}cclll@{}}
\caption{MAD Goat Kubernetes Vulnerability Inventory}\label{tab:vuln-inventory}\\
\toprule
\textbf{ID} & \textbf{Title} & \textbf{2025} & \textbf{2022} & \textbf{CWE} & \textbf{CIS} & \textbf{Scope} \\
\midrule
\endfirsthead
\multicolumn{7}{c}{\tablename\ \thetable{} (continued)}\\
\toprule
\textbf{ID} & \textbf{Title} & \textbf{2025} & \textbf{2022} & \textbf{CWE} & \textbf{CIS} & \textbf{Scope} \\
\midrule
\endhead
\midrule
\multicolumn{7}{r}{\textit{Continued on next page}}\\
\endfoot
\bottomrule
\endlastfoot
\texttt{k01\-runAsRoot} & Run-as-root containers & K01 & K01 & CWE-250 & 5.2.6 & chart \\
\texttt{k01\-readOnlyRootFs} & Writable root filesystem & K01 & K01 & CWE-276 & --- & chart \\
\texttt{k01\-privileged} & Privileged containers & K01 & K01 & CWE-250 & 5.2.2 & chart \\
\texttt{k01\-unboundedResources} & Unbounded resource usage & K01 & K01 & CWE-400 & --- & chart \\
\texttt{k02\-clusterAdminBinding} & Default SA bound to cluster-admin & K02 & K03 & CWE-269 & 5.1.1 & chart \\
\texttt{k02\-secretsListWatch} & Role with list/watch on Secrets & K02 & K03 & CWE-732 & 5.1.3 & chart \\
\texttt{k03\-plaintextConfigMapSecrets} & Credentials in plaintext ConfigMap & K03 & K08 & CWE-522 & 5.4.1 & chart \\
\texttt{k03\-jwtKeyInEnv} & JWT signing key in env (ConfigMap) & K03 & K08 & CWE-522 & 5.4.2 & chart \\
\texttt{k04\-noPSA} & Namespace without Pod Security Admission & K04 & K04 & CWE-693 & 1.2.9 & cluster \\
\texttt{k04\-weakKyverno} & Kyverno installed with no-op rules & K04 & K04 & CWE-693 & 1.2.12 & cluster \\
\texttt{k05\-noDefaultDeny} & No default-deny NetworkPolicy & K05 & K07 & CWE-923 & 5.3.2 & chart \\
\texttt{k05\-permissiveNetpol} & Overly permissive label-based NetworkPolicy & K05 & K07 & CWE-923 & 5.3.1 & chart \\
\texttt{k05\-noEgressControls} & Missing egress controls & K05 & K07 & CWE-923 & 5.3.1 & chart \\
\texttt{k06\-traefikDashboardExposed} & Traefik dashboard publicly exposed & K06 & --- & CWE-284 & --- & chart \\
\texttt{k06\-rabbitmqMgmtExposed} & RabbitMQ management UI publicly exposed & K06 & --- & CWE-284 & --- & chart \\
\texttt{k07\-kubeletAnonymousAuth} & kubelet anonymous-auth=true & K07 & K09 & CWE-306 & 1.2.1 & cluster \\
\texttt{k07\-kubeletAlwaysAllow} & kubelet authorization-mode=AlwaysAllow & K07 & K09 & CWE-306 & 1.2.6 & cluster \\
\texttt{k09\-anonymousKeycloakBootstrap} & Keycloak default admin credentials + start-dev & K09 & K06 & CWE-200 & --- & chart \\
\texttt{k09\-sharedServiceAccount} & All pods share the default ServiceAccount & K09 & K06 & CWE-200 & --- & chart \\
\texttt{k10\-noAuditPolicy} & No apiserver audit policy configured & K10 & K05 & CWE-778 & 1.2.16 & cluster \\
\end{longtable}


\subsection{Vulnerability Catalogue}
\label{sec:k8s-catalogue}